\chapter{Deep Learning dalam Praktik}
\label{ch:praktik}

Teori deep learning saya dapat di semester pertama. Cara melatih jaringan dengan benar, dan cara mencari tahu mengapa
sebuah jaringan tidak mau belajar, lebih banyak datang dari latihan-latihan kuliah Deep Learning and Neural Networks II.
Enam tugasnya bergerak dari dasar PyTorch, jaringan visi dan transfer learning, pemantauan
eksperimen dan CNN lanjutan, autoencoder dan VAE, contoh adversarial dan GAN, sampai normalizing flows. Arsitektur-arsitekturnya
sudah dibahas di \cref{ch:arsitektur}. Bab ini mengumpulkan sisi lainnya: apa yang terjadi di dalam kerangka kerja ketika kita
menulis satu baris \texttt{loss.backward()}, dan kebiasaan kecil yang memisahkan eksperimen yang berhasil dari yang membuang
waktu berhari-hari.

\section{Graf komputasi dan autograd}

Kuliah membuka dengan sejarah singkat kerangka kerja deep learning. Generasi awal seperti Theano dan TensorFlow 1 memakai graf
statis: seluruh komputasi didefinisikan dulu sebagai graf simbolik, lalu dijalankan berkali-kali. PyTorch \citep{paszke2019}
memakai graf dinamis, yang dibangun ulang setiap kali kode dijalankan (\emph{define-by-run}), sehingga kontrol alur Python biasa,
seperti percabangan dan perulangan, bisa dipakai di dalam model. TensorFlow 2 mengambil jalan serupa dengan \emph{gradient tape}
yang merekam operasi selama lintasan maju.

Di balik layar, setiap tensor yang dihasilkan operasi menyimpan rujukan ke operasi yang menghasilkannya dan ke input-input
operasi itu. Ketika \texttt{backward()} dipanggil pada loss skalar, kerangka kerja menelusuri graf ini mundur dan menerapkan
aturan rantai, persis reverse-mode autodiff di \cref{ch:dlbasic}. Tiga perilaku sering membingungkan pemula, dan ketiganya
langsung masuk akal dari gambaran ini. Gradien \emph{diakumulasikan} ke atribut \texttt{.grad}, tidak ditimpa, sehingga
\texttt{optimizer.zero\_grad()} harus dipanggil sebelum setiap lintasan mundur. Kalau lupa, langkah kedua memakai jumlah gradien
langkah pertama dan kedua. Memanggil \texttt{.detach()} memotong graf, sehingga tidak ada gradien yang mengalir melewati tensor
itu, berguna misalnya untuk target network di DQN (\cref{ch:rl}). Dan blok \texttt{torch.no\_grad()} menghentikan perekaman graf
sama sekali, sehingga evaluasi tidak memakan memori untuk menyimpan aktivasi yang tidak akan diturunkan.

Satu putaran pelatihan, dengan urutan yang menjadi bahan soal di kuliah, berbentuk seperti ini:
\begin{verbatim}
model.train()
for x, y in loader:
    x, y = x.to(device), y.to(device)
    optimizer.zero_grad()      # bersihkan gradien lama
    loss = loss_fn(model(x), y) # lintasan maju, graf direkam
    loss.backward()            # lintasan mundur, gradien diisi
    optimizer.step()           # perbarui parameter
\end{verbatim}
Untuk evaluasi, \texttt{model.eval()} mematikan dropout dan membuat batch normalization memakai statistik bergerak, dan
\texttt{torch.no\_grad()} membungkus lintasan maju. Data dibaca oleh \texttt{Dataset}, yang tahu cara mengambil satu contoh, dan
\texttt{DataLoader}, yang menyusunnya menjadi mini-batch, mengacak urutannya, dan memuatnya secara paralel dengan beberapa
proses pekerja agar GPU tidak menunggu.

\section{Resep men-debug jaringan}

Jaringan yang tidak belajar jarang memberi pesan galat. Ia hanya menghasilkan kurva loss yang datar atau angka yang tidak masuk
akal. Beberapa pemeriksaan murah menangkap sebagian besar masalah.

Pemeriksaan pertama adalah loss awal. Untuk klasifikasi $K$ kelas dengan softmax dan cross-entropy, jaringan yang baru
diinisialisasi seharusnya memberi peluang hampir seragam, $1/K$ untuk setiap kelas, sehingga loss awalnya mendekati
$-\ln(1/K)=\ln K$. Untuk sepuluh kelas, itu sekitar $2{,}30$. Kalau loss awal jauh lebih besar, misalnya 30, logit awalnya terlalu
besar dan jaringan sudah yakin pada jawaban yang salah sebelum belajar apa pun, tanda inisialisasi atau normalisasi input
yang keliru.

Pemeriksaan kedua adalah \emph{overfit} satu batch kecil. Ambil beberapa contoh saja dan latih berulang-ulang pada contoh yang
sama. Jaringan dengan kapasitas apa pun seharusnya bisa menghafalnya sampai loss mendekati nol. Kalau tidak bisa, masalahnya
ada di kode, misalnya label yang tidak sesuai dengan input, gradien yang tidak mengalir karena \texttt{detach} yang salah tempat,
atau laju belajar yang keliru, bukan di jumlah data atau arsitektur.

Pemeriksaan ketiga adalah laju belajar. Naikkan laju belajar secara eksponensial dari kecil ke amat besar selama
beberapa ratus langkah, dan catat loss di setiap langkah. Loss mula-mula hampir tidak berubah, lalu turun cepat, lalu meledak.
Laju belajar yang baik biasanya sedikit di bawah titik loss turun paling curam, jauh sebelum ledakan. Kurva ini menghubungkan
praktik dengan teori di \cref{ch:optimasi}: titik ledakan adalah tempat langkah melampaui batas $2/L$.

Pemeriksaan terakhir adalah melihat apa yang dilihat jaringan. \citet{zeiler2014} memproyeksikan aktivasi lapisan dalam kembali
ke ruang piksel dengan jaringan dekonvolusi, memakai indeks lokasi maksimum dari pooling untuk ``membuka'' pooling dan
transposed convolution untuk membalik konvolusi. Lapisan awal ternyata merespons tepi dan warna, lapisan tengah tekstur, dan
lapisan akhir bagian objek. Tugas kuliah meminta kami mereplikasi teknik ini dengan hanya mempertahankan beberapa aktivasi
terbesar di setiap peta fitur. Gambar seperti ini sering langsung memperlihatkan bahwa jaringan memperhatikan latar belakang
alih-alih objek, sama dengan pelajaran Grad-CAM di \cref{ch:xai}.

\section{Memori dan presisi}

Ketika model atau batch tidak muat di memori GPU, dua teknik sering dipakai. \istilah{Gradient accumulation} memecah satu batch
besar menjadi beberapa batch kecil, menjalankan lintasan maju dan mundur untuk masing-masing tanpa memanggil
\texttt{optimizer.step()}, lalu memperbarui parameter sekali di akhir. Karena gradien diakumulasikan, hasilnya sama dengan satu
batch besar, asal loss setiap batch kecil dibagi dengan jumlah batch kecil. Batch 256 bisa dijalankan sebagai empat batch 64,
dengan setiap loss dikalikan $1/4$. Pengecualiannya adalah batch normalization, yang tetap melihat statistik batch kecil.

\istilah{Mixed precision} \citep{micikevicius2018} menghitung sebagian besar operasi dalam bilangan floating point 16 bit, yang
memakan separuh memori dan jauh lebih cepat di GPU modern, sambil menyimpan salinan utama bobot dalam 32 bit. Masalahnya ada
pada rentang angka. Format setengah presisi standar hanya bisa mewakili bilangan normal dari sekitar $6\times10^{-5}$ sampai
$65\,504$. Gradien yang lebih kecil dari batas bawah, yang umum di jaringan dalam, menjadi nol. Obatnya adalah \istilah{loss
scaling}: kalikan loss dengan faktor besar, misalnya 1024, sebelum lintasan mundur, sehingga semua gradien ikut membesar
1024 kali dan keluar dari daerah yang hilang, lalu bagi kembali sebelum memperbarui bobot. Faktornya bisa disesuaikan otomatis:
diperkecil kalau muncul angka tak hingga, diperbesar perlahan kalau tidak.

\section{Mencari hyperparameter}

Laju belajar, ukuran batch, weight decay, jumlah lapisan, dan peluang dropout harus dipilih sebelum pelatihan. Kuliah membahas
beberapa strategi. Pencarian manual mengandalkan intuisi. Pencarian grid mencoba semua kombinasi dari beberapa nilai untuk
setiap hyperparameter. Pencarian acak menarik kombinasi secara acak.

\citet{bergstra2012} memberi argumen sederhana mengapa pencarian acak biasanya lebih baik dari grid. Misalkan dari dua
hyperparameter hanya satu yang benar-benar berpengaruh, keadaan yang ternyata umum. Grid $3\times3$ menghabiskan sembilan
percobaan tetapi hanya mencoba tiga nilai berbeda dari hyperparameter yang penting. Sembilan percobaan acak mencoba sembilan
nilai berbeda (\cref{fig:gridrandom}). Ada juga argumen peluang yang tidak bergantung pada dimensi. Kalau kita ingin
setidaknya satu percobaan jatuh di lima persen konfigurasi terbaik, peluang $n$ percobaan acak semuanya gagal adalah
$0{,}95^n$. Dengan 60 percobaan, peluang berhasil $1-0{,}95^{60}\approx95$ persen, berapa pun jumlah hyperparameter-nya.

\begin{figure}[ht]
\centering
\begin{tikzpicture}[x=1cm,y=1cm]
  \begin{scope}
    \draw[abu] (0,0) rectangle (3,3);
    \foreach \i in {0.5,1.5,2.5}{\foreach \j in {0.5,1.5,2.5}{\fill[biru] (\i,\j) circle (2.2pt);}}
    \foreach \i in {0.5,1.5,2.5}{\draw[oranye,thick] (\i,-0.1) -- (\i,-0.3);}
    \node[label] at (1.5,3.35) {grid};
    \node[label] at (1.5,-0.6) {hyperparameter penting};
  \end{scope}
  \begin{scope}[xshift=4.5cm]
    \draw[abu] (0,0) rectangle (3,3);
    \foreach \i/\j in {0.3/2.1,0.8/0.6,1.1/2.7,1.4/1.3,1.9/0.3,2.2/2.3,2.5/1.0,2.8/1.8,0.55/1.6}{\fill[biru] (\i,\j) circle (2.2pt); \draw[oranye,thick] (\i,-0.1) -- (\i,-0.3);}
    \node[label] at (1.5,3.35) {acak};
    \node[label] at (1.5,-0.6) {hyperparameter penting};
  \end{scope}
\end{tikzpicture}
\caption{Sembilan percobaan dengan grid dan dengan pencarian acak. Kalau hanya sumbu mendatar yang berpengaruh, grid hanya
menguji tiga nilainya, sedangkan pencarian acak menguji sembilan.}
\label{fig:gridrandom}
\end{figure}

\istilah{Bayesian optimization} \citep{snoek2012} memakai hasil percobaan sebelumnya untuk memilih percobaan berikutnya. Sebuah
model pengganti, biasanya Gaussian process dari \cref{ch:mlat}, memodelkan skor validasi sebagai fungsi hyperparameter, lengkap
dengan ketidakpastiannya. Titik berikutnya dipilih dengan memaksimalkan sebuah fungsi akuisisi. Yang paling umum adalah
\istilah{expected improvement}: kalau skor terbaik sejauh ini $f^*$ dan model meramal skor di titik $\vx$ berdistribusi
$\N(\mu,\sigma^2)$, perbaikan harapan atas $f^*$ adalah
\begin{equation}
\mathrm{EI}(\vx)=(\mu-f^*)\,\Phi(z)+\sigma\,\varphi(z),\qquad z=\frac{\mu-f^*}{\sigma},
\end{equation}
dengan $\Phi$ dan $\varphi$ fungsi distribusi dan kerapatan normal baku. Suku pertama menghargai titik yang diramal bagus,
suku kedua menghargai titik yang tidak pasti, keseimbangan eksplorasi dan eksploitasi yang sama dengan bandit di
\cref{ch:klasik}. Dua titik dengan $\mu$ sama dengan $f^*$ punya $z=0$, sehingga $\mathrm{EI}=\sigma\varphi(0)\approx0{,}40\,\sigma$:
titik yang lebih tidak pasti lebih menarik untuk dicoba. Neural architecture search memperluas gagasan ini ke struktur
jaringan itu sendiri, dengan harga komputasi yang jauh lebih besar.

Untuk semua strategi, pemantauan eksperimen menjadi kebutuhan, bukan kemewahan. Kuliah memakai alat pencatat eksperimen
seperti Weights \& Biases: setiap percobaan mencatat hyperparameter, kurva loss latih dan validasi, dan versi kode. Tanpa catatan
seperti itu, eksperimen terbaik minggu lalu tidak bisa diulang.

\section{Transfer learning lebih dekat}

\Cref{ch:arsitektur} menyebut resep dasar transfer learning: bekukan lapisan yang sudah dilatih, latih ulang kepalanya. Di
PyTorch, membekukan berarti mematikan \texttt{requires\_grad} pada parameter yang tidak ingin diubah, sehingga autograd tidak
menghitung gradiennya dan optimizer tidak menyentuhnya. Gambar input juga harus disesuaikan dengan apa yang dilihat jaringan
saat pretraining, misalnya diperbesar ke $224\times224$ dan dinormalisasi dengan rata-rata dan simpangan baku ImageNet.

Lapisan mana yang sebaiknya dibekukan? \citet{yosinski2014} mengukur seberapa mudah fitur di setiap lapisan dipindahkan dari satu
tugas ke tugas lain. Fitur lapisan awal bersifat umum dan hampir selalu bisa dipakai ulang. Fitur lapisan akhir semakin
spesifik pada tugas asal. Mereka juga menemukan bahwa memindahkan fitur lalu menyetel ulang seluruh jaringan memberi
generalisasi yang lebih baik daripada melatih dari nol, bahkan ketika data tugas baru cukup banyak. Dalam praktik, semakin
sedikit data dan semakin mirip tugasnya, semakin banyak lapisan yang dibekukan. Semakin banyak data dan semakin berbeda
tugasnya, semakin banyak lapisan yang disetel ulang, biasanya dengan laju belajar yang lebih kecil untuk lapisan awal.

\section{Mengulang hasil}

Hasil deep learning sulit diulang persis. Bobot awal, urutan data, dropout, dan augmentasi semuanya acak, dan beberapa operasi
GPU tidak deterministik karena urutan penjumlahan floating point yang berubah-ubah memberi hasil yang sedikit berbeda.
Menetapkan benih acak untuk semua sumber keacakan membuat satu proses bisa diulang. Tetapi benih yang berbeda tetap memberi hasil
yang berbeda, kadang cukup jauh. Karena itu kebiasaan yang jujur adalah menjalankan beberapa benih dan melaporkan rata-rata
beserta sebarannya. Selisih dua metode yang lebih kecil dari sebaran antar-benih belum bisa disebut perbaikan.

\section{Di luar kelas}

Kebiasaan di bab ini berlaku di semua bab penerapan. Pemeriksaan loss awal dan overfit satu batch sama bergunanya untuk PINN
(\cref{ch:pinn}) dan simulator graf (\cref{ch:operator}) seperti untuk klasifikasi gambar. Mixed precision memungkinkan model
fondasi cuaca dan protein dilatih dalam waktu yang wajar. Bayesian optimization dipakai bukan hanya untuk hyperparameter jaringan,
tetapi juga untuk mengkalibrasi parameter simulasi yang mahal, di mana setiap evaluasi berarti satu simulasi penuh. Dan laporan
dengan beberapa benih dan pita sebaran adalah pengingat yang sama dengan error bar di \cref{ch:stat}: angka tanpa
ketidakpastian hanya separuh cerita.

\sumberbab{Bab ini mengikuti tugas dan materi kuliah Deep Learning and Neural Networks II (PyTorch dan autograd, jaringan visi
dan transfer learning, pemantauan eksperimen dan pencarian hyperparameter, CNN lanjutan, transposed convolution dan
visualisasi, autoencoder dan VAE, contoh adversarial dan GAN, normalizing flows). Makalah pendukung: \citet{paszke2019},
\citet{bergstra2012}, \citet{snoek2012}, \citet{yosinski2014}, \citet{zeiler2014}, dan \citet{micikevicius2018}.}
